\section{Both resonances and every logarithmic moment}\label{bmh:sec:resonances}

\subsection{An all-weight double-zeta evaluation}
Write
\begin{equation}\label{bmh:eq:Idef}
 I_{m,n}^{(h)}=\int_0^\infty
 \frac{\log^h x\,\Li_m(-x)\Li_n(-x)}{x(1+x)}\dd x,
 \qquad I_{m,n}=I_{m,n}^{(0)}.
\end{equation}

\begin{theorem}[Compact double-zeta reduction]\label{bmh:thm:compact}
For every $m,n\ge1$, with $W=m+n$,
\begin{align}
 \boxed{\begin{aligned}
 I_{m,n}={}&\left(\binom W{m-1}+\binom W{n-1}\right)\zeta(W+1)\\
 &+\sum_{j=1}^{W-1}j\,A_j(m,n)\zeta(j+1,W-j).
 \end{aligned}}
 \label{bmh:eq:compact}
\end{align}
where
\begin{equation}\label{bmh:eq:Aj}
 A_j(m,n)=\binom{W-j-1}{m-1}+\binom{W-j-1}{n-1}
           -\binom{j-1}{m-1}-\binom{j-1}{n-1}.
\end{equation}
In particular $A_{W-j}(m,n)=-A_j(m,n)$. All double-zeta values in the
formula are ordinarily convergent and have weight $W+1$.
\end{theorem}

\begin{proof}[A direct proof at resonance]
The $a=0$ limit of~\eqref{bmh:eq:resolvent} is
\[
 I_0(c,d)=\frac{c\log c/(c-1)-d\log d/(d-1)}{c-d}.
\]
On the cone $u=v+t$ it is the following difference of two nonnegative
kernels:
\begin{equation}\label{bmh:eq:resonant-cone}
 I_0(e^{v+t},e^v)=
 \frac{t e^{-(v+t)}}{(1-e^{-t})(1-e^{-(v+t)})}
 -\frac{v e^{-(2v+t)}}{(1-e^{-(v+t)})(1-e^{-v})}.
\end{equation}
Each term, after multiplication by $(v+t)^{m-1}v^{n-1}$, is integrable.
This can be seen either from its elementary behavior at the axes or from
the convergent double sums below. Consequently Tonelli's theorem applies
separately to their geometric expansions.

Using the same binomial expansion as in the proof of the master, the
first cone contributes
\begin{align}
 \sum_{k=0}^{m-1}\binom{W-2-k}{n-1}
 \bigl[&(k+1)\{\zeta(k+2,W-1-k)+\zeta(W+1)\}\notag\\
       &-(W-1-k)\zeta(W-k,k+1)\bigr].
 \label{bmh:eq:resonant-unsimplified}
\end{align}
The weak inequality in the first geometric sum accounts for the
single-zeta diagonal. The second cone gives the same expression with
$m,n$ interchanged. Reindexing the positive and negative double-zeta
terms gives $jA_j(m,n)$. The coefficient of the diagonal is
\[
 \sum_{k=0}^{m-1}(k+1)\binom{W-2-k}{n-1}
 +\sum_{k=0}^{n-1}(k+1)\binom{W-2-k}{m-1}
 =\binom W{m-1}+\binom W{n-1},
\]
a double application of the hockey-stick identity. This proves the
formula independently of any numerical recognition.
\end{proof}

\subsection{All Taylor coefficients without singular coordinates}
Let
\begin{equation}\label{bmh:eq:even-coeff}
 \frac{\pi u}{\sin\pi u}=\sum_{j\ge0}e_{2j}u^{2j},
 \qquad e_0=1,\quad
 e_{2j}=2(1-2^{1-2j})\zeta(2j)\quad(j\ge1).
\end{equation}
For positive integers $A,B$ and $\nu\ge1$, define
\begin{align}
 c_\nu(A,B)={}&-\frac{(A)_\nu}{\nu!}\zeta(A+\nu,B)\notag\\
 &+\sum_{v=1}^{\nu}
 \frac{(A)_{\nu-v}(B)_v}{(\nu-v)!v!}
 \bigl\{\zeta(B+v,A+\nu-v)+\zeta(A+B+\nu)\bigr\}.
 \label{bmh:eq:cjet}
\end{align}
Every first double-zeta index in this expression is at least two. An
equivalent, simpler corner formula is
\begin{equation}\label{bmh:eq:c11jet}
 c_\nu(1,1)=\frac12\sum_{v=1}^{\nu-1}
       \zeta(v+1)\zeta(\nu-v+1)
       +\frac{\nu+1}{2}\zeta(\nu+2).
\end{equation}
Define $T_\nu(m,n)$ by replacing every $\bmhC_{A,B}$ in
\eqref{bmh:eq:K} by $c_\nu(A,B)$.

\begin{lemma}[Cancelled coefficient identity]\label{bmh:lem:cjet}
For $|u|<1$,
\begin{equation}\label{bmh:eq:C-Taylor}
 \bmhC_{A,B}(1-u)=\sum_{\nu\ge1}c_\nu(A,B)u^\nu.
\end{equation}
\end{lemma}
\begin{proof}
For $A\ge2$, expand the normally convergent double Hurwitz series and
use
\[
 D_B(1-u)=\sum_{v\ge1}\frac{(B)_v}{v!}\zeta(B+v)u^v.
\]
The latter also holds at $B=1$ by the digamma expansion. At total degree
$\nu$, the terms of $Z_{1-u}(A,B)$ with positive shift in the second
index cancel against the corresponding stuffle term of
$\zeta(A;1-u)D_B(1-u)$. What remains is exactly
\eqref{bmh:eq:cjet}. For $A=1$, take the limit from spectral $A>1$ in
the analytic combined cone integral of~\eqref{bmh:eq:cone-moment}.
All coordinates surviving on the right of~\eqref{bmh:eq:cjet} converge
at this limit. The same argument reaches $A=B=1$; alternatively,
expanding~\eqref{bmh:eq:Ccorner} directly gives
\eqref{bmh:eq:c11jet}. Local uniform convergence permits coefficient
extraction throughout.
\end{proof}

\begin{theorem}[Complete resonant logarithmic moments]
\label{bmh:thm:jets}
For every $m,n\ge1$ and $h\ge0$,
\begin{equation}\label{bmh:eq:all-jets}
 \boxed{\displaystyle
 I_{m,n}^{(h)}=h!\sum_{j=0}^{\lfloor h/2\rfloor}
                      e_{2j}T_{h+1-2j}(m,n).}
\end{equation}
Every term has representation weight $m+n+h+1$. At $h=0$ the formula
reduces to Theorem~\ref{bmh:thm:compact}.
\end{theorem}
\begin{proof}
Insert~\eqref{bmh:eq:C-Taylor} into~\eqref{bmh:eq:master} and use
\eqref{bmh:eq:even-coeff}. The constant term of $\bmhK(1-u)$ is zero, so
its factor of $u$ cancels the cosecant pole. The coefficient of $u^h$
is the right side of~\eqref{bmh:eq:all-jets} divided by $h!$.
Differentiating the original integral is justified by
Proposition~\ref{bmh:prop:strip}. The weight assertion follows by
adding the indices in~\eqref{bmh:eq:cjet} and the even-zeta weight.
\end{proof}

\begin{remark}[Depth is not the same as the generating algebra]
For $h=0,1$ the answer is a linear combination of single and double zeta
values. For larger $h$, the factor $e_{2j}$ may multiply a double value.
It is therefore correct to say that the answer belongs to the algebra
generated by ordinary and double zeta values. It is \emph{not} correct
to call every such product additive depth two. Expanding by stuffle gives
additive depth at most three. No minimal-depth assertion is made.
\end{remark}

\subsection{The second resonance at \texorpdfstring{$a=-1$}{a=-1}}
Define
\begin{equation}\label{bmh:eq:minus-coeff}
 c^-_\nu(A,B)=c_\nu(A,B)
 -\sum_{v=1}^{\nu}
    \frac{(A)_{\nu-v}(B)_v}{(\nu-v)!v!}\zeta(B+v),
\end{equation}
and let $T^-_\nu(m,n)$ be the corresponding replacement in
\eqref{bmh:eq:K}.

\begin{corollary}[Negative Mellin resonance]\label{bmh:cor:negative}
For $m,n\ge1$ and $h\ge0$,
\begin{equation}\label{bmh:eq:negative-jets}
 \boxed{\displaystyle
 \int_0^\infty\frac{\log^h x\,\Li_m(-x)\Li_n(-x)}{x^2(1+x)}\dd x
 =-h!\sum_{j=0}^{\lfloor h/2\rfloor}e_{2j}T^-_{h+1-2j}(m,n).}
\end{equation}
\end{corollary}
\begin{proof}
The shift law gives
\[
 \bmhC_{A,B}(2-u)=\bmhC_{A,B}(1-u)-(1-u)^{-A}D_B(1-u).
\]
Multiplying the two elementary Taylor series proves that its coefficients
are $c^-_\nu$. Now use $a=-1+u$ and
$\sin\pi(-1+u)=-\sin\pi u$ in~\eqref{bmh:eq:master}.
\end{proof}
For instance,
\begin{equation}\label{bmh:eq:negative11}
 \int_0^\infty\frac{\log^2(1+x)}{x^2(1+x)}\dd x
 =2\zeta(2)-2\zeta(3).
\end{equation}
This resonance exists because the product is $O(x^2)$ at zero. The
one-factor lower endpoint $a>-1$ cannot be reused for the bilinear
problem.

\subsection{Small-weight reductions with explicit algebra}
The first values are
\begin{align}
 I_{1,1}&=2\zeta(3),\label{bmh:eq:example11}\\
 I_{2,1}&=\frac{17}{4}\zeta(4),\label{bmh:eq:example21}\\
 I_{3,1}&=\frac{11}{2}\zeta(5)+\zeta(2)\zeta(3),\label{bmh:eq:example31}\\
 I_{2,2}&=2\zeta(5)+4\zeta(2)\zeta(3).\label{bmh:eq:example22}
\end{align}
Here is enough exact algebra to verify the reductions. The depth-two sum
formula
\begin{equation}\label{bmh:eq:sumformula}
 \sum_{k=2}^{w-1}\zeta(k,w-k)=\zeta(w),\qquad w\ge3,
\end{equation}
follows without regularization: sum the finite geometric progression
inside $\sum_{n>m}$ to obtain
\[
 \sum_{k=2}^{w-1}\frac1{n^k m^{w-k}}
 =\frac{1/(n m^{w-2})-1/n^{w-1}}{n-m}.
\]
The first part summed over $n>m$ is
$\sum_m H_m/m^{w-1}$, and the second is
$\sum_n H_{n-1}/n^{w-1}$. Their difference is $\zeta(w)$.
At weight four, stuffle gives $\zeta(2,2)=3\zeta(4)/4$, so
$\zeta(3,1)=\zeta(4)/4$. At weight five, use the convergent shuffle
and stuffle rows
\begin{align*}
 \zeta(2)\zeta(3)&=\zeta(2,3)+3\zeta(3,2)+6\zeta(4,1),\\
 \zeta(2)\zeta(3)&=\zeta(2,3)+\zeta(3,2)+\zeta(5),
\end{align*}
together with~\eqref{bmh:eq:sumformula}. Solving gives
\[
 \zeta(4,1)=2\zeta(5)-\zeta(2)\zeta(3),\quad
 \zeta(2,3)=\frac92\zeta(5)-2\zeta(2)\zeta(3).
\]
Substitution into~\eqref{bmh:eq:compact} proves the displayed examples.
